\documentclass[11pt]{article}
\usepackage{times}
\pagestyle{empty}
\usepackage[override]{cmtt}

\usepackage[pdftex]{hyperref}
\hypersetup{
	pdftitle={Curriculum Vitae},
	pdfauthor={Michael Hicks}, 
	pdfproducer={PDFlatex},
	pdfkeywords={},
	colorlinks=true,
	citecolor=black,
	linkcolor=black,
	urlcolor=blue,
	bookmarksopen=true,
        pdfpagemode=UseNone
}

\newcommand{\rate}[1]{
 {\em #1}
}

\topmargin -0.5in \headheight\baselineskip
\textheight 9.25in              % 1in top and bottom margin
\textwidth 6.5in        % 1in left and right margin
\oddsidemargin 0in
\evensidemargin 0in

\setcounter{secnumdepth}{0}
\makeatletter
\def\section#1{\@startsection{section}{1}{\z@}{1.2em plus -.25em minus -.25em}{.3em plus .1em}{\normalsize\bf}{#1}}
\makeatother

\parindent 0in
\parskip 0em

\begin{document}

\newcommand{\papertitle}[1]{\textbf{#1}}

\begin{center}
{\LARGE\bf Michael W. Hicks}
\\
April 28, 2026\\
\end{center}

\begin{description}
\setlength{\labelsep}{0ex}
\setlength{\itemsep}{-0.5ex}

\item {\bf{\textsf Contact Information}}

\item

  \begin{tabular}{p{.45\textwidth}p{.45\textwidth}} 
Email: \url{mwh@cis.upenn.edu} 
& LinkedIn: \href{https://www.linkedin.com/in/michael-hicks-a053311/}{\texttt{michael-hicks-a053311}}\\
WWW: \url{https://mhicks.me}
    & Blog: \url{http://pl-enthusiast.net} \\
\end{tabular}

\end{description}

\begin{description}
\setlength{\labelsep}{0ex}
\setlength{\itemsep}{-0.5ex}

\item {\bf\textsf Education} 

\item BS, with honors, Pennsylvania State University 1993 (Computer
  Science).  
\item MS, University of Pennsylvania 1996 (Computer and
  Information Science).
\item PhD, University of Pennsylvania August, 2001 (Computer and
  Information Science).

\end{description}

\begin{description}
\setlength{\labelsep}{0ex}
\setlength{\itemsep}{-0.5ex}

\item {\bf\textsf Appointments} 

 
\item {\em Cecilia Fitler Moore Professor} (August 2025--present), \\ Dept of Computer and Information Science, University of Pennsylvania.
\item {\em Associate Chair} (January 2026--present), \\ Dept of Computer and Information Science, University of Pennsylvania.
\item {\em Director} (August 2025--present), Schlein Center for Cybersecurity, University of Pennsylvania.
\item {\em Amazon Scholar} (July 2025--present), Amazon Web Services.
\item {\em Senior Principal Applied Scientist} (January 2022--July 2025), Amazon Web Services
\item {\em Professor Emeritus} (January 2023--present), Computer Science, University of Maryland.
\item {\em Chief Technology Officer} (September 2018--December 2021), Correct Computation, Inc.
\item {\em Professor} (July 2013--December 2022), Computer Science, University of Maryland.\\ (Assistant and Associate professor from July 2002-2013.)
\item {\em Associate Co-Chair for Undergraduate Education} (September 2017--December 2021), \\ Department of Computer Science, University of Maryland.
\item {\em Consultant} (August--November 2020) Covington and Burling, LLP.
\item {\em Consultant} (June 2019--July 2020), Stealth Software Technologies, Inc.
\item {\em Consultant and testifying witness} (February 2014–2017), Wilson Sonsini Goodrich and Rosati.
\item {\em Visiting Researcher} (June 2015--Aug. 2015), Microsoft
  Research, Redmond.
\item {\em Director} (October 2011--2013), Maryland Cybersecurity
  Center, University of Maryland.
\item {\em Affiliate Faculty} (2005--2011), Electrical
  Engineering, University of Maryland.
\item {\em Adjunct} (2006--2015), Institute for Defense Analyses
   Center for Computing Sciences, Bowie, MD.
\item {\em Visiting Researcher} (Sep. 2008--Nov. 2008), Microsoft
  Research, Cambridge, the United Kingdom.
\item {\em Post-doctoral Research Associate} (2001--2002), Computer Science,
  Cornell University, Ithaca, New York. 
\item {\em Scientist} (1998--1999), NEC Research Institute, Princeton,
  New Jersey. 
\item {\em Software Engineer} (1993--1996), ARINC, Inc.
\end{description}

\begin{description}
\setlength{\labelsep}{0ex}
\setlength{\itemsep}{-0.5ex}


        \item {\bf\textsf Notable awards and service}

\item \textbf{Fellow of the ACM}, named January 2023.
\item \textbf{ACM SIGPLAN Chair}, served 2015--2018; for the next 3
  years I served as \textbf{Past Chair}. SIGPLAN is
  one of 30 special interests groups (SIGs) of ACM members, in this
  case with a focus on programming languages (PLAN). SIGPLAN hosts
  five high-profile conferences among a slate of 13. There are nearly
  3,000 SIGPLAN members.
\item \textbf{ACM SIGPLAN Distinguished Service Award}, named June
  2022, for contributions as SIGPLAN Chair and Past Chair, and for service on
  program committees and other SIGPLAN activities over the years.
\item \textbf{Editor in Chief}, \emph{Proceedings of the ACM on Programming Languages (PACMPL)}, November 2022--2025. PACMPL is the journal that publishes the proceedings of the ``big four'' SIGPLAN conferences (POPL, PLDI, OOPSLA, ICFP).
\item \textbf{University of Maryland Distinguished Scholar-Teacher},
  2015--2016. This award is given to about five University of Maryland
  faculty each year, from among the 1,500 across all Departments on
  campus, who demonstrate remarkable achievement in both research and
  teaching.
 \item \textbf{Department of Computer Science Faculty Teaching Award}, 2003,
   2005, 2007, 2014.
 \item \textbf{NSF CAREER Award}, 2004.
\item \textbf{Program committee (PC) Chair and member of high-profile
    conferences}. First Program Chair of IEEE Cybersecurity
  Development Conference (SecDev), which focuses on building systems
  security, held November 2016; served on the SC and PC 2017--2019. Co-founder, Workshop on Programming Languages for Quantum Computing (PLanQC), 2020. 
  Other PC memberships include IEEE
  Computer Security Foundations Symposium (2014, 2025, Chair 2015 \& 2016), IEEE
  Symposium on Security and Privacy (PC member 2009, 2015, 2017,
  2022, 2026, Area Chair 2018), ACM Conference on Computer and Communications Security (PC
  member, 2008 and 2017, Area Chair 2019), ACM SIGPLAN Symposium on
  Principles of Programming Languages (PC member 2009 \& 2024, Chair 2012, SC
  Chair 2018--2021), ACM SIGPLAN Conference on Programming Language
  Design and Implementation (PC member, 2007, 2013, 2016, Area Chair
  2023), ACM SIGPLAN SPLASH/OOPSLA Conference (PC member 2007, 2011, 2014, 2021, 2026).
 \item \textbf{Associate Editor}, Transactions on Programming Languages and
   Systems (TOPLAS), January 2012--2017.
 \item \textbf{Member, National Academy of Sciences committee} on
   Enhancing the Assurance and Nimbleness of Large-Scale Integrated
   Software-Based Systems, November 2023--August 2025.
\item \textbf{Member, DARPA Information Science and Technology Board (ISAT)},
  September 2009--2012, September 2026--forward.  This panel consists of 30
    individuals (largely faculty and researchers) from
    U.S. institutions which advises DARPA by conducting studies on
    topics of increasing relevance and impact. I co-led a study that
    spurred the creation of a DARPA program (HACMS) and boosted
    synthesis research (mentioned as a "hot topic" at PLDI’22).  
% \item PC Member, ACM SIGPLAN Workshop on Types in
%   Language Design and Implementation (TLDI), January 2011.
% \item PC Member, ACM SIGPLAN International
%   Conference on Functional Programming (ICFP), August 2010.
% \item Review Committee (RC) Member, ACM SIGPLAN Conference on
%   Programming Language Design and Implementation (PLDI), June 2010.
% \item Tutorials Chair, ACM SIGPLAN Conference on
%   Programming Language Design and Implementation (PLDI), June 2010.
% \item PC Member, ACM SIGPLAN-SIGSOFT Workshop on
%   Program Analysis for Software Tools and Engineering (PASTE), June 2010.
% \item PC Member, European Symposium on Programming (ESOP), March 2010.
\item \textbf{Member, CyberMaryland 2.0 Task Force}, April--October
  2012, Panel commissioned by the Maryland Department of Business and
  Economic Development (DBED) to explore how to increase Maryland's
  impact and standing in the cybersecurity arena.
 \item \textbf{Member, University System of Maryland Chancellor’s Task Force on Cybersecurity}, Nov 2010–Feb 2011.
% \item $\quad\bullet$ Participant, National Science Foundation, Computer
%   Systems Research Future Directions Workshop, March 2010.
% \item PC Member, ACM Symposium on the Principles of Programming
%   Languages (POPL), 2009.
% \item $\quad\bullet$ 
% \item Member of DARPA Information Processing Technology Office (IPTO)
% Futures Panel, Fall 2007.  {\scriptsize This panel consists of a few select
% faculty gathered to help IPTO chart its research goals and to develop
% potential programs.}
%\item \textbf{Member of DARPA Computer Science Study Panel} (CS2P),
 % 2006--2010, which consists of a dozen select junior faculty from
 % U.S. institutions, convening four times during 2006, and twice %   during 2007.}
% \item PC Member, ACM SIGPLAN Conference on
%   Object-Oriented Programming Languages, Systems, and Applications
%   (OOPSLA), 2007.
% \item PC Member, ACM SIGPLAN Conference on
%   Programming Language Design and Implementation (PLDI), 2007.
% \item $\quad\bullet$ General Chair and Technical Program Chair, Workshop on Programming
%   Languages and Analysis for Security (PLAS), 2007.
% \item PC Member, Workshop on Programming
%   Languages and Analysis for Security (PLAS), 2006.
% \item Developed methodology for improving management of academic
%   research groups; adopted by dozens of computer scientists and
%   others; cf. Michael Hicks and Jeffrey~S. Foster, Score: Agile
%   research group management, {\em Communications of the {ACM}},
%   53(10):30--31, October 2010.
\item \textbf{Founder and member, SIGPLAN Committee on Empirical
    Evaluations}. I founded the committee while SIGPLAN Chair,
  inspired by what I found with my fuzzing work. We produced a
  checklist of best practices at
  \url{http://sigplan.org/Resources/EmpiricalEvaluation/} 
\item \textbf{Founder and first editor in chief, PL Perspectives}, the
  blog of the ACM Special Interest Group on Programming Languages
  (SIGPLAN), curating more than 90 articles. \url{https://blog.sigplan.org}
  % \item \textbf{Co-conceived and organized \emph{Build-it, Break-it,
%       Fix-it}}, a security-oriented programming and exploitation contest.
%  \url{https://builditbreakit.org}. 
% \item Co-taught \emph{CMSC 498B: Secure Maryland}, Spring 2012, 
%   on penetration testing.  Students could attack
%   University of Maryland networks; they found, and helped fix,
%   hundreds of
%   vulnerabilities. Cf. \url{http://www.cs.umd.edu/~mwh/securemd/CMSC498B__Secure_Maryland/Home.html} 
\end{description}

\begin{description}
\setlength{\labelsep}{0ex}
\setlength{\itemsep}{-0.5ex}


\item {\bf\textsf Scientific publications}

\item Co-author of 153 peer-reviewed scientific papers
\item 14750 citations and H-index of 58, per Google Scholar\\
  (\url{https://scholar.google.com/citations?user=Gggzp7UAAAAJ&hl=en})
 \item All papers at \url{https://mhicks.me/papers/mwh.html}

\item
\item {\it\textsf Papers in last 5 years, and other notables}
\item (2025)

\item
Aleks Chakarov, Jaco Geldenhuys, Matthew Heck, Michael Hicks, Sam Huang,
  Georges-Axel Jaloyan, Anjali Joshi, K.~Rustan~M. Leino, Mikael Mayer, Sean
  McLaughlin, Akhilesh Mritunjai, Clement Pit-Claudel, Sorawee Porncharoenwase,
  Marianna Rapoport, Cody Roux, Neha Rungta, Robin Salkeld, Matthias
  Schlaipfer, Daniel Schoepe, Johanna Schwartzentruber, Serdar Tasiran, Aaron
  Tomb, Jean-Baptiste Tristan, Emina Torlak, Lucas Wagner, Michael~W. Whalen,
  Remy Willems, Tongtong Xiang, Tae~Joon Byun, Joshua Cohen, Ruijie Fang,
  Junyoung Jang, Jakob Rath, Hira~Taqdees Syeda, Dominik Wagner, and Yongwei
  Yuan.
\papertitle{Formally verified cloud-scale authorization.}
In {\em Proceedings of the International Conference on Software
  Engineering (ICSE)}, April 2025. \textbf{Distinguished Paper}.

\item
\item (2024)

\item Joseph~W. Cutler, Craig Disselkoen, Aaron Eline, Shaobo He, Kyle Headley,
  Michael Hicks, Kesha Hietala, Eleftherios Ioannidis, John Kastner, Anwar
  Mamat, Darin McAdams, Matt McCutchen, Neha Rungta, Emina Torlak, and Andrew
  Wells.
\papertitle{Cedar: A new language for expressive, fast, safe, and analyzable
  authorization.}
In {\em Proceedings of the {ACM} Conference on Object-Oriented
  Programming Languages, Systems, and Applications (OOPSLA)}, April 2024.
  
\item Craig Disselkoen, Aaron Eline, Shaobo He, Kyle Headley, Michael Hicks, Kesha
  Hietala, John Kastner, Anwar Mamat, Matt McCutchen, Neha Rungta, Bhakti Shah,
  Emina Torlak, and Andrew Wells.
\papertitle{How we built Cedar: A verification-guided approach.}
In {\em Proceedings of the European Software Engineering Conference
  and the ACM SIGSOFT Symposium on the Foundations of Software Engineering
  (ESEC/FSE)}, July 2024.
Industry papers track.

% \item
% Joseph~W. Cutler, Michael Hicks, and Emina Torlak.
% newblock Improving the stability of type safety proofs in Dafny, January 2024.
% newblock Dafny Workshop talk.

\item
\item (2023)

\item
  Kesha Hietala, Robert Rand, Liyi Li, Shih-Han Hung, Xiaodi Wu, and Michael
  Hicks.
\papertitle{A verified optimizer for quantum circuits.}
{\em {ACM} Transactions on Programming Languages and Systems
  (TOPLAS)}, July 2023.
Extends POPL'21 paper.
  
\item
Yuxiang Peng, Kesha Hietala, Runzhou Tao, Liyi Li, Robert Rand, Michael Hicks,
  and Xiaodi Wu.
\papertitle{A formally certified end-to-end implementation of Shor's
  factorization algorithm.}
{\em Proceedings of the National Academy of Sciences}, 120(21), May
  2023.
  
\item
Jie Zhou, John Criswell, and Michael Hicks.
\papertitle{ Fat pointers for temporal memory safety of C.}
 In {\em Proceedings of the {ACM} Conference on Object-Oriented
  Programming Languages, Systems, and Applications (OOPSLA)}, April 2023.

\item
Haowei Deng, Yuxiang Peng, Michael Hicks, and Xiaodi Wu.
\papertitle{Automating {NISQ} application design with meta quantum circuits with
  constraints ({MQCC}).}
{\em {ACM} Transactions on Quantum Computing (TQC)}, 4(3), April 2023.

\item
Zenong Zhang, George Klees, Eric Wang, Michael Hicks, and Shiyi Wei.
\papertitle{Fuzzing configurations of program options.}
{\em {ACM} Transactions on Software Engineering and Methodology
  (TOSEM)}, 32(2), March 2023.
A shorter version appeared at the 2022 Fuzzing Workshop.

\item
Michael Coblenz, April Porter, Varun Das, Teja Nallagorla, and Michael Hicks.
\papertitle{A multimodal study of challenges using {Rust}.}
In {\em Proceedings of the Workshop on the Evaluation and Usability
  of Programming Languages and Tools (PLATEAU)}, February 2023.

\item
Ian Sweet, David Darais, David Heath, Ryan Estes, Bill Harris, and Michael
  Hicks.
\papertitle{ Symphony: Expressive secure multiparty computation with coordination.}
 {\em $\langle$Programming$\rangle$}, 7(14), February 2023.


\item
Finn Voichick, Robert Rand, and Michael Hicks.
\papertitle{ {Qunity}: A unified language for quantum and classical computing.}
 In {\em Proceedings of the {ACM} Conference on Principles of
  Programming Languages (POPL)}, January 2023.

  \item
  \item (2022)

\item
Liyi Li, Finnegan Voichick, Kesha Hietala, Yuxiang Peng, Xiaodi Wu, and Michael
  Hicks.
\papertitle{ Verified compilation of quantum oracles.}
 In {\em Proceedings of the {ACM} Conference on Object-Oriented
  Programming Languages, Systems, and Applications (OOPSLA)}, December 2022.

\item
Aravind Machiry, John Kastner, Matt McCutchen, Aaron Eline, Kyle Headley, and
  Michael Hicks.
\papertitle{ {C} to {Checked C} by {3C}.}
 In {\em Proceedings of the {ACM} Conference on Object-Oriented
  Programming Languages, Systems, and Applications (OOPSLA)}, December 2022.
 \textbf{Distinguished Paper}.
   
\item
Kelsey~R. Fulton, Daniel Votipka, Desiree Abrokwa, Michelle~L. Mazurek, Michael
  Hicks, and James Parker.
\papertitle{ Understanding the how and the why: Exploring secure development
  practices through a course competition.}
 In {\em Proceedings of the {ACM} Conference on Computer and
  Communications Security (CCS)}, October 2022.

\item
Liyi Li, Yiyun Liu, Deena~L. Postol, Leonidas Lampropoulos, David~Van Horn, and
  Michael Hicks.
\papertitle{ A formal model of {Checked C}.}
 In {\em Proceedings of the Computer Security Foundations Symposium
  (CSF)}, August 2022.

\item
Zenong Zhang, Zach Patterson, Michael Hicks, and Shiyi Wei.
\papertitle{ Fixreverter: A realistic bug injection methodology for benchmarking
  fuzz testing.}
 In {\em Proceedings of the USENIX Security Symposium (USENIX SEC)},
  August 2022.
 \textbf{Distinguished Paper}.

\item
Michael Coblenz, Michelle Mazurek, and Michael Hicks.
\papertitle{ Does the bronze garbage collector make rust easier to use? a
  controlled experiment.}
 In {\em Proceedings of the International Conference on Software
  Engineering (ICSE)}, May 2022.
\textbf{Distinguished Paper Award Nominee}.

\item
Zenong Zhang, George Klees, Eric Wang, Michael Hicks, and Shiyi Wei.
\papertitle{ Fuzzing configurations of program options.}
 In {\em Proceedings of the International Fuzzing Workshop}, April
  2022.

%   \item (2021)

% \item
% Kelsey Fulton, Anna Chan, Dan Votipka, Michael Hicks, and Michelle Mazurek.
% \papertitle{ Benefits and drawbacks of adopting a secure programming language:
%   {Rust} as a case study.}
%  In {\em Proceedings of the Symposium on Usable Privacy and Security
%   (SOUPS)}, August 2021.

% \item
% Souvik Bhattacherjee, Gang Liao, Michael Hicks, and Daniel~J. Abadi.
% \papertitle{ Bullfrog: Online schema evolution via lazy evaluation.}
%  In {\em Proceedings of the {ACM} SIGMOD International Conference on
%   Management of Data (SIGMOD)}, June 2021.

% \item
% Kesha Hietala, Liyi Li, Akshaj Gaur, Aaron Green, Robert Rand, Xiaodi Wu, and
%   Michael Hicks.
% \papertitle{ Applying and expanding the {VOQC} toolkit.}
%  In {\em Informal Proceedings of the Workshop on Programming Languages
%   and Quantum Computing (PLanQC)}, June 2021.

% \item
% Kesha Hietala, Robert Rand, Shih-Han Hung, Liyi Li, and Michael Hicks.
% \papertitle{ Proving quantum programs correct.}
%  In {\em Proceedings of the Conference on Interative Theorem Proving
%   (ITP)}, June 2021.

% \item
% Ian Sweet, David Darais, David Heath, Ryan Estes, William Harris, and Michael
%   Hicks.
% \papertitle{ Symphony: A concise language model for {MPC}.}
%  In {\em Informal Proceedings of the Workshop on Foundations on
%   Computer Secuirty (FCS)}, June 2021.

% \item
% Finn Voichick and Michael Hicks.
% \papertitle{ Toward a quantum programming language for higher-level formal
%   verification.}
%  In {\em Informal Proceedings of the Workshop on Programming Languages
%   and Quantum Computing (PLanQC)}, June 2021.

\item
\item (Other notable publications)

  % \item
  % ThanhVu Nguyen, Timos Antopoulos, Andrew Ruef, and
  % Michael Hicks. A Counterexample-guided Approach to Finding Numerical
  % Invariants. In {\em Proceedings of the European Software Engineering
  % Conference and the ACM SIGSOFT Symposium on the Foundations of
  % Software Engineering (ESEC/FSE)}, September 2017.

% \item Karla Saur, Michael
%   Hicks, and Jeffrey S. Foster. C-Strider: Type-Aware Heap Traversal
%   for C. \emph{Software, Practice, and Experience, }
%   May 2015.
  
% \item
% Chang Liu, Yan Huang, Elaine Shi, Jonathan Katz, and Michael Hicks.
%  Automating efficient {RAM}-model secure computation.
%  In {\em Proc. of the {IEEE} Symposium on Security and Privacy
%   (Oakland)}, May 2014.

% \item Andrew Miller, Michael Hicks, Jonathan Katz, and Elaine
%   Shi.  Authenticated data structures, generically.  In {\em
%     Proc. of the {ACM} Conference on Principles of Programming
%     Languages (POPL)}, January 2014.

% \item Christopher M. Hayden, Edward K. Smith, Eric A. Hardisty,
%   Michael Hicks, and Jeffrey S. Foster. Evaluating Dynamic Software
%   Update Safety Using Efficient Systematic Testing. {\em IEEE
%     Transactions on Software Engineering}, 38(6):1340-1354, December
%   2012. 

% \item
% Trevor Jim, Nikhil Swamy, and Michael Hicks.
%  Defeating scripting attacks with browser-enforced embedded policies.
%  In {\em Proc. of the International World Wide Web Conference
%   (WWW)}, May 2007.

% \item
% Brian~J. Corcoran, Nikhil Swamy, and Michael Hicks.
% Cross-tier, label-based security enforcement for web applications.
% In {\em Proceedings of the {ACM} SIGMOD International Conference on
%   Management of Data (SIGMOD)}, pages 269--282, June 2009.

\item
Kesha Hietala, Robert Rand, Shih-Han Hung, Xiaodi Wu, and Michael Hicks.
\papertitle{ A verified optimizer for quantum circuits.}
 In {\em Proceedings of the {ACM} Conference on Principles of
  Programming Languages (POPL)}, January 2021.
\textbf{Distinguished Paper}.


\item 
  Daniel Votipka, Kelsey Fulton, James Parker, Matthew Hou, Michelle~L. Mazurek,
  and Michael Hicks.
\papertitle{Understanding security mistakes developers make: Qualitative analysis
  from {Build It, Break It, Fix It}.}
In {\em Proceedings of the USENIX Security Symposium (USENIX SEC)},
  August 2020. \textbf{Distinguished Paper}.


\item George T. Klees, Andrew Ruef, Benjamin Cooper, Shiyi Wei, and
  Michael Hicks.
  \papertitle{Evaluating Fuzz Testing.}
  In \emph{Proceedings of the ACM
  Conference on Computer and Communications Security (CCS)}, October
  2018. \textbf{NSA Best Scientific Cybersecurity Paper 2018}.

%  \item
% Timos Antonopoulos, Paul Gazzillo, Michael Hicks, Eric Koskinen, Tachio
%   Terauchi, and Shiyi Wei.
%  Decomposition instead of self-composition for proving the absence of
%   timing channels.
%  In {\em Proceedings of the {ACM} Conference on Programming Language
%   Design and Implementation (PLDI)}, June 2017.

    
\item
  Chang Liu, Austin Harris, Martin Maas, Michael Hicks, Mohit Tiwari,
  and Elaine Shi.
\papertitle{  Ghostrider: A hardware-software system for memory trace oblivious computation.}
  In {\em
    Proc. of the International Conference on Architectural
    Support for Programming Languages and Operating Systems (ASPLOS)},
  March 2015.   \textbf{Best Paper}.
 
\item
Christopher~M. Hayden, Karla Saur, Edward~K. Smith, Michael Hicks, and
  Jeffrey~S. Foster.
\papertitle{Efficient, general-purpose dynamic software updating for {C}.}
{\em {ACM} Transactions on Programming Languages and Systems
  (TOPLAS)}, 36(4):13, October 2014.
  

 \item
 Aseem Rastogi, Matthew~A. Hammer, and Michael Hicks.
\papertitle{  Wysteria: A programming language for generic, mixed-mode multiparty computations.}
  In {\em Proc. of the {IEEE} Symposium on Security and Privacy
   (Oakland)}, May 2014.

% \item Matthew Hammer, Yit Phang Khoo, Michael Hicks, and Jeffrey
%   S. Foster. Adapton: Composable, Demand-Driven Incremental
%   Computation.  In {\em Proceedings of the ACM Conference on
%     Programming Language Design and Implementation (PLDI)}, June 2014.
  
\item Chang Liu, Michael Hicks, and Elaine Shi.
\papertitle{  Memory trace oblivious program execution.}
   In {\em Proc. of
     the Computer Security Foundations Symposium (CSF)}, June
   2013. \textbf{NSA Best Scientific Cybersecurity Paper 2013}.

% Piotr Mardziel, Mario Alvim, Michael Hicks, and Michael Clarkson.
%  Quantifying information flow for dynamic secrets.
%  In {\em Proc. of the {IEEE} Symposium on Security and Privacy
%   (Oakland)}, May 2014.

\item  
  Yit Phang Khoo, Jeffrey S. Foster, and Michael Hicks.
\papertitle{  Expositor: Scriptable Time-Travel Debugging with First Class Traces.}
  In
  \emph{Proceedings of the International Conference on Software Engineering
  (ICSE)}, May 2013.

% \item
% Lu\'is Pina, Lu\'is Veiga, and Michael Hicks.
%  Rubah: {DSU} for java on a stock {JVM}.
%  In {\em Proc. of the {ACM} Conference on Object-Oriented
%   Programming Languages, Systems, and Applications (OOPSLA)}, October 2014.

% \item
% Matthew Hammer, Yit~Phang Khoo, Michael Hicks, and Jeffrey~S. Foster.
%  Adapton: Composable, demand-driven incremental computation.
%  In {\em Proc. of the {ACM} Conference on Programming Language
%   Design and Implementation (PLDI)}, June 2014.

% \item
% Piotr Mardziel, Stephen Magill, Michael Hicks, and Mudhakar Srivatsa.
%  Dynamic enforcement of knowledge-based security policies using
%   probabilistic abstract interpretation.
%  {\em Journal of Computer Security}, 21:463--532, October 2013.

% \item Jong hoon (David)~An, Avik Chaudhuri, Jeffrey~S. Foster, and
%   Michael Hicks.  Dynamic inference of static types for Ruby.  In {\em
%     Proc. of the {ACM} Conference on Principles of Programming
%     Languages (POPL)}, January 2011.

% \item
% Aseem Rastogi, Piotr Mardziel, Matthew Hammer, and Michael Hicks.
%  Knowledge inference for optimizing secure multi-party computation.
%  In {\em Proc. of the {ACM SIGPLAN} Workshop on Programming
%   Languages and Analysis for Security (PLAS)}, June 2013.

% \item
%   Jong hoon (David) An,
%   Avik Chaudhuri, Jeffrey S. Foster, and Michael Hicks.   Dynamic
%   Inference of Static Types for Ruby.  In {\em Proceedings
%   of the ACM Conference on Principles of Programming Languages (POPL)},
%   pages 459-472, January 2011.
  
% \item
%   Jean-Philippe Martin, Michael Hicks, Manuel Costa, Periklis Akritidis, and
%    Miguel Castro.
%  Dynamically checking ownership policies in concurrent {C/C++}
%    programs.
%  In {\em Proceedings of the {ACM} Conference on Principles of
%    Programming Languages (POPL)}, January 2010.


% \item Piotr Mardziel, Michael Hicks, Jonathan Katz, and Mudhakar Srivatsa.
% Knowledge-oriented secure multiparty computation.
% {\em Proc. of the {ACM SIGPLAN} Workshop on Programming
%   Languages and Analysis for Security (PLAS)}, June 2012.

     
\item Kin-Keung Ma, Yit Phang Khoo,
  Jeffrey S. Foster, and Michael Hicks.
  \papertitle{Directed Symbolic Execution.}
  In Eran Yahav, editor,
  {\em Proceedings of the Static Analysis Symposium (SAS)}, volume 6887 of
  {\em Lecture Notes in Computer Science}, pages 95-111. Springer, September
  2011.
    
%  \item
%  Yit~Phang Khoo, Jeffrey~S. Foster, Michael Hicks, and Vibha Sazawal.
%   Path projection for user-centered static analysis tools.
%   In {\em Proceedings of the {ACM} Workshop on Program Analysis for
%    Software Tools and Engineering (PASTE)}, November 2008.
   
 \item Saurabh Srivastava, Michael Hicks, Jeffrey~S. Foster, and
   Patrick Jenkins.
\papertitle{   Modular Information Hiding and Type Safe Linking for {C}.}
   {\em IEEE Transactions on Software Engineering (TSE)},
   34(3):1-20, May 2008.

% \item Iulian Neamtiu, Michael Hicks, Jeffrey~S. Foster, and Polyvios
%   Pratikakis.  Contextual effects for version-consistent dynamic
%   software updating and safe concurrent programming.  In {\em
%     Proc. of the {ACM} Conference on Principles of Programming
%     Languages (POPL)}, January 2008.

 \item Nick~L. Petroni, Jr. and Michael Hicks.
   \papertitle{Automated detection of persistent kernel control-flow attacks.}
   In {\em Proc. of the
     {ACM} Conference on Computer and Communications Security (CCS)},
   October 2007. 

% \item Polyvios Pratikakis, Jeffrey~S. Foster, and Michael Hicks.
%   Context-sensitive correlation analysis for detecting races.  In {\em
%     Proc. of the {ACM} Conference on Programming Language Design
%     and Implementation (PLDI)}, pages 320--331, June 2006.

% \item Iulian Neamtiu, Michael Hicks, Gareth Stoyle, and Manuel Oriol.
%   Practical dynamic software updating for {C}.  In {\em Proc. of
%     the {ACM} Conference on Programming Language Design and
%     Implementation (PLDI)}, pages 72--83, June 2006.

   \item
   Nikhil Swamy, Michael Hicks, Greg Morrisett, Dan Grossman, and Trevor Jim.
\papertitle{ Safe manual memory management in {Cyclone}.}
 {\em Science of Computer Programming (SCP)}, 62(2):122--144, October
  2006.
 Special issue on memory management. \emph{Key ideas 
   now embodied in the Rust programming language.}

   
\item Trevor Jim, Greg Morrisett, Dan Grossman, Michael Hicks, James
  Cheney, and Yanling Wang.
  \papertitle{{C}yclone: A safe dialect of {C}.}
  In
  {\em Proc. of the {USENIX} Annual Technical Conference}, 
  June 2002.

 \item Dan Grossman, Greg Morrisett, Trevor Jim, Michael Hicks, Yanling
   Wang, and James Cheney.
\papertitle{   Region-based memory management in {C}yclone.}
   In {\em Proc. of the {ACM} Conference on
     Programming Language Design and Implementation (PLDI)}, pages
   282--293. {ACM}, June 2002.

\item 
Michael Hicks.
 {\em Dynamic Software Updating}.
 PhD thesis, Department of Computer and Information Science,
  University of Pennsylvania, August 2001.
\textbf{ACM SIGPLAN Doctoral Dissertation Award~2002.}

\end{description}

\begin{description}
\setlength{\labelsep}{0ex}
\setlength{\itemsep}{-0.5ex}


\item {\bf\textsf Patents}

\item Emina Torlak, Darin McAdams, Neha Rungta, Michael W. Hicks, Craig Ryan Disselkoen, Aaron Joseph Eline, John Holman Kastner, Kyle Headley, Anwar Mamat, Richar Matthew McCutchen, Andrew Marshall Wells, Kesha Hanne Hietala, Shaobo He, Mark Edward Stalzer, and Julian Lovelock.
  \papertitle{Authorization Policy Evaluation}.
  US12483595B2, United States Patent and Trademark Office, Nov. 25, 2025.

\item Emina Torlak, Kyle Headley, Michael W. Hicks, Neha Rungta, Andrew Marshall Wells. 
  \papertitle{Authorization Policy Analysis}. 
  US12425455B2, United States Patent and Trademark Office, Sep. 23, 2025. 

\item Michael W. Hicks, John Holman Kastner, Emina Torlak, Richard
  Matthew McCutchen, Darin McAdams, Neha Rungta, Aaron Joseph Eline,
  Joseph Wallace Cutler, Eleftherios
  Ioannidis. \papertitle{Authorization policy
    validation}. US20240179182A1, United States Patent and Trademark Office, March 25,
  2025. 

\item Michael W. Hicks, B. S. Lavender, Anwar Mamat, and Ian Nicholas
  Sweet. \papertitle{Automated generation of source code models from
    machine code for code analysis}.
  US11321218B1, United States Patent and Trademark Office, May 5,
  2022.

\end{description}

\begin{description}
\setlength{\labelsep}{0ex}
\setlength{\itemsep}{-0.5ex}

\item {\bf\textsf Notable invited talks}

\item \textit{Measuring End-User Harm.} Invited talk, 17th Cyber Public Health workshop, Nov 2025.
\item \textit{Cedar: A language for expressing fast, safe, and
    fine-grained authorization policies}
        \begin{itemize}
            \item Keynote, IEEE Secure Development Conference, Oct 2025.
            \item UPenn CIS Colloquium, Nov 2024.
            \item Invited talk, Cylab Seminar, Carnegie Mellon University, Feb 2024.
            \item Distinguished lecture, Dept of CS, Iowa State
              University, Nov 2023.
         \end{itemize}

    \item \textit{Amazon's Formal Methods Journey.} Invited talk, DARPA Resilient Software Systems Colloquium, June 2025.
    \item
      \textit{Contesting Secure Software Development}
      \begin{itemize}
        \item Invited speaker, ISAT/DARPA RAPIDS Workshop, March 2023.
        \item Keynote, Computer Security Foundations Symposium, June
          2020.
         \item UPenn CIS Colloquium, Nov 2019.
          \end{itemize}
\item 
        \textit{From Verified Compilation to Shor's Algorithm.} Keynote, Second International Workshop on Programming Languages for Quantum Computing (PLanQC 2021), June 2021.
    \item
        \textit{Evaluating Fuzz Testing}
        \begin{itemize}
            \item Keynote, Meeting on Hot Topics on the Science of Security (HotSOS), Sep 2020.
            \item
                Best scientific cybersecurity paper award acceptance talk, National Security Agency, Oct 2019.
            \item
                Invited lecture, 9th International Summer School on Information Security and Protection, Canberra, Australia, July 2018. 
            \item
                Keynote, Cornell Programming Languages Retreat, Sep 2018. 
             \item
                Colloquium talk, EPFL Department of Computer and Communications Sciences, Lausanne, Switzerland, Oct 2018.
        \end{itemize}
    \item
        \textit{Building Security In: Programming Language-based Techniques for Ensuring Software Security.} Invited lecture, 9th International Summer School on Information Security and Protection, Canberra, Australia, July 2018.
    \item
        \textit{Languages for Oblivious Computation.} Keynote talk at the ACM Workshop on Programming Languages and Analyses for Security (PLAS), co-located with the ACM SIGSAC Conference on Computer and Communications Security (CCS), Oct 2017. 
    \item
        \textit{Programming Languages Meet Cryptography.} Series of three invited lectures at the Cornell, Maryland, Max Planck Pre-doctoral Research School, August 8-13, 2017 in Saarbruecken, Germany. 
    \item
        \textit{Tackling Software Insecurity with Abstract Interpretation.} Invited talk, Workshop on the Next 40 Years of Abstract Interpretation (N40AI), collocated with POPL, Jan 2017.
   \item
        \textit{Authenticated Data Structures, Generically.} Colloquium talk, Max Planck Institute for Software Systems (MPI-SWS), Saarbrucken, Germany, Dec 2014. 
    \item
        \textit{Memory-Trace Oblivious Program Execution.} NSA Best Scientific Cybersecurity Award Paper presentation, Sep 2014.
    \item
        \textit{On-line patching for better security, and other security challenges.} Colloquium talk, Johns Hopkins Applied Physics Lab, May 2013. 
    \item
        \textit{Expositor: Scriptable time-travel debugging with first class traces.} Colloquium talk, Max Planck Institute for Software Systems (MPI-SWS), Saarbrucken, Germany, Nov 2012.
    \item
        \textit{Kitsune: Efficient, General-purpose Dynamic Software Updating for C.} Colloquium talk, Johns Hopkins University, Department of Computer Science, Nov 2012.
    \item
        \textit{Cybersecurity: Past, Present, Looking ahead.} Keynote, Maryland Francis King Carey School of Law Symposium on \textit{Cybersecurity: Safeguarding Information in a Digital Age}, March 2012.
    \item
       \textit{Dynamic Enforcement of Knowledge-based Security Policies.} Colloquium talk, George Washington University, March 2012
    \item
        \textit{Dynamic inference of static types for Ruby.} Colloquium talk, Department of Computer Science, Cornell University, Oct 2011
    \item
        \textit{Software Synthesis for cost-effective development of correct, efficient software.} ISAT outbrief to DARPA program managers and the deputy director on October 4, 2011. I co-chaired this ISAT study with Armando Solar-Lezama of MIT. 
    \item
        \textit{Practical Dynamic Software Updating for C}
        \begin{itemize}
            \item
                Colloquium talk, University of Kent at Canterbury Computer Laboratory, Mar 2009.
            \item
                Colloquium talk, University of Cambridge Computer Laboratory, Oct 2008.
            \item
                Colloquium talk, Department of Computer Science, the University of Virginia, Dec 2006. 
            \item
                Colloquium talk, Department of Computer Science, the Pennsylvania State University, Nov 2005. 
        \end{itemize}
    \item
        \textit{SELINKS: A language for provably secure web applications.} Distinguished Lecturer talk, IBM T.J. Watson Research Laboratory, July 2007. 
    \item
        \textit{LOCKSMITH: Context-sensitive Correlation Analysis for Detecting Races.} Colloquium talk, Department of Computer Science, University of Massachusetts (Amherst), Nov 2006. 
    \item
        \textit{Static Analysis to Improve Software Quality} (with Jeff Foster), National Research Council of Canada, Institute for Information Technology invited talk, June 2006. 
    \item
        \textit{Safe and Flexible Memory Management in Cyclone.} Colloquium talk, Department of Computer Science, Purdue University, Aug 2003. 
    \item
        \textit{Dynamic Software Updating.} Colloquium talk, The Computer Laboratory at Cambridge University, Cambridge, England, Oct 2001. 

        \end{description}

        

\begin{description}
\setlength{\labelsep}{0ex}
\setlength{\itemsep}{-0.5ex}

\item {\bf\textsf PhDs and postdoctoral researcher advisees} 
\item I have supervised 13 postdoctoral researchers, and
  advised, co-advised, or served on the committees of 65 graduate
  students at UMD and elsewhere. PhD students and postdocs I
  closely and personally mentored are given below.
\item
  
\item Liyi Li, Postdoc 2020--2023 (Assistant Professor, Iowa State University)
\item Ethan Cecchetti, Postdoc 2021--2022 (Assistant Professor, University of Wisconsin)
\item Michael Coblenz, Postdoc 2020--2022 (Assistant Professor, UCSD)
\item Kesha Hietala, PhD 2022 (Researcher, Sandia National Labs)
\item Ian Sweet, PhD 2022 (Software Research Engineer, Galois)
\item Robert Rand, Postdoc 2018--2020 (Assistant Professor, University
  of Chicago)
\item Leonidas Lampropoulos, Postdoc 2018--2020 (Assistant Professor,
  University of Maryland)  
\item James Parker, PhD 2020 (Software Research Engineer, Galois)
\item Andrew Ruef, PhD 2018 (First job: Quantitative Researcher IDA/CCS)
\item Chang Liu, PhD 2016 (Researcher, Citadel Securities)
\item Shiyi Wei, Postdoc 2015--2017 (Associate Professor, University
  of Texas at Dallas)
\item Aseem Rastogi, PhD 2016 (Principal Researcher, Microsoft Research)
\item Luis Pina, PhD 2015 (Assistant Professor, University of Illinois
  Chicago)
\item Matthew Hammer, Postdoc 2013--2015 (First job: Assistant Professor, CU Boulder)
\item Karla Saur, PhD 2015 (Distributed Systems Engineer, Nvidia DGX Cloud)
\item Piotr Mardziel, PhD 2015 (First job: Systems Scientist, Carnegie Mellon)
\item Khoo Yit Phang, PhD 2013 (Software Engineering Manager, Mathworks)
\item Nataliya Guts, Postdoc 2011--2013 (Security Solutions Engineer, Futurae)
\item Chris Hayden, PhD 2012 (Senior Software Engineer, Amazon Web Services)
\item Stephen Magill, Postdoc 2010--2012 (VP, Product Innovation at Sonatype)
\item Justin McCann, PhD 2012 (Principal Software Engineering Manager, Microsoft)
\item Kin Keung (Martin) Ma, PhD 2011 (Software Engineer, Google)
\item Saurabh Srivastava, PhD 2010 (CEO and Founder, Consequent AI)
\item Pavlos Papageorge, PhD 2008 (Software Engineer, Google (AI))
\item Nikhil Swamy, PhD 2008 (Senior Principal Researcher, Microsoft)
\item Iulian Neamtiu, PhD 2008 (Professor, NJIT)
\item Polyvios Pratikakis, PhD 2008 (Associate Professor, University of Crete, Greece)
\item Nick L. Petroni, Jr., PhD 2008 (Chief Scientist, Volexity)
\item Manuel Oriol, Postdoc 2006--2008 (Professor, Constructor Institute of Technology, Switzerland)
\end{description}

\begin{description}
\setlength{\labelsep}{0ex}
\setlength{\itemsep}{-0.5ex}

\item {\bf\textsf Other notable mentoring activities}

\item Undergraduate advising -- mentored 36 UMD undergraduate students
  on research projects. Notables:
  \begin{itemize}
    \item Yiyun Liu (2018-2020) 2020 CRA outstanding
      undergraduate researcher, honorable mention. Now PhD student at
      UPenn.
    \item Andrew Ruef (2010-2013) and James Parker (2012-2013)
      later got PhDs with me.
    \item Ted Smith (2009-2013), co-authored 4 peer-reviewed papers on
      dynamic software updating.
      \end{itemize}
\item High-school student project advising -- mentored 9 students on
  summer projects. Notables:
  \begin{itemize}
  \item Akshaj Gaur (2020) co-authored workshop paper \emph{Applying
      and Expanding the VOQC Toolkit}, now an MS student at UMD.
    \item George Klees (2017) first-authored paper \emph{Evaluating
        Fuzz Testing} that won the NSA's Best Scientific Cybersecurity
      Paper competition in 2018.
    \end{itemize}

  \end{description}

\begin{description}
\setlength{\labelsep}{0ex}
\setlength{\itemsep}{-0.5ex}

\item {\bf\textsf Educational activities}

\item CIS 7000-010, Empirical Security \& Privacy, for Humans, 10 students, Fall 2025.
\item CIS 7010-003, Secure Software Engineering \& Management, 14 students, Spring 2026.
\item (\underline{Courses taught at UMD})
  \item CMSC 631, Program Analysis and Understanding, 30 students, Fall 2021.
    \item CMSC 330, Organization of Programming Languages, 145 students, Spring 2021.
    \item CMSC 330, Organization of Programming Languages, 145 students, Spring 2019.
    \item CMSC 330, Organization of Programming Languages, 145 students, Spring 2018.
    \item CMSC 631, Program Analysis and Understanding, 28 students, Fall 2017.
    \item CMSC 330, Organization of Programming Languages, 215 students, Spring 2017.
    \item CMSC 396H, Honors Seminar, 13 students, Fall 2016.
    \item CMSC 838G, Mechanized Proof and Verified Software, 8 students, Spring 2016.
    \item CMSC 330, Organization of Programming Languages, 129 students, Fall 2015.
    \item CMSC 330, Organization of Programming Languages, 130 students, Spring 2015.
    \item CMSC 433, Programming Language Technologies and Paradigms, 50 students, Fall 2014
    \item CMSC 838G, Software Security, 10 students, Spring 2014
    \item CMSC 433, Programming Language Technologies and Paradigms, 50 students, Fall 2013
    \item CMSC 631, Program Analysis and Understanding, 26 students, Spring 2013
    \item CMSC 330, Organization of Programming Languages, 113 students, Spring 2013
    \item CMSC 498L, Cybersecurity lab (co-taught with 4 other profs), 25 students, Fall 2012
    \item CMSC 498B, Secure Maryland (penetration testing), 7 students, Spring 2012
    \item CMSC 631, Program Analysis and Understanding, 13 students, Fall 2011
    \item CMSC 838G, Software Security, 5 students, Spring 2011
    \item CMSC 433, Programming Language Technologies and Paradigms, 40 students, Fall 2010
    \item CMSC 330, Organization of Programming Languages, 66 students, Spring 2010
    \item CMSC 631, Program Analysis and Understanding, 9 students, Fall 2009
    \item CMSC 631, Program Analysis and Understanding, 16 students, Fall 2007
    \item CMSC 412, Operating Systems, 42 students, Spring 2007
    \item CMSC 631, Program Analysis and Understanding, 20 students, Fall 2006
    \item CMSC 433, Programming Language Technologies and Paradigms, 44 students, Spring 2006
    \item CMSC 412, Operating Systems, 50 students, Fall 2005
    \item CMSC 838Z, Language-based Security, 7 students, Spring 2005
    \item CMSC 412, Operating Systems, 45 students, Fall 2004
    \item CMSC 838Z, Tools and Techniques for Software Dependability, 8 students, Spring 2004
    \item CMSC 433, Programming Language Technologies and Paradigms, 48 students, Fall 2003
    \item CMSC 838Y, Agile and Adaptive Programming, 10 students, Spring 2003
    \item CMSC 433, Programming Language Technologies and Paradigms, 47 students, Fall 2002
    \item (Other teaching experiences)
    \item \textit{Software Security} - I developed a massively-open
    on-line course (MOOC) under the UMD-Coursera agreement, on this
    topic, as part of a specialization on cybersecurity. I fully
    developed the content for this course, including videos, lecture
    notes, and projects. See
    \url{https://mhicks.me/software_security_course/index.html}. When
    retired in December 2023, this course had had 70,000+ participants.  
\item \textit{Cybersecurity Capstone} – I developed the Capstone
  project for the UMD specialization on cybersecurity, based on our
  Build-it, Break-it, Fix-it contest. We hosted two capstones in 2015
  and one in 2016. See \url{https://builditbreakit.org}
      
\end{description}

\begin{description}
\setlength{\labelsep}{0ex}
\setlength{\itemsep}{-0.5ex}

\item {\bf\textsf Grants, contracts, and gifts}

\item Total of 34 grants awarded and industry gifts received, for a
  total of \$15.875M raised. 

\item

\item AFOSR FA95502110051 Software Assurance for Quantum Programs, \$218,875, January 2021-2024. PI: Michael Hicks. (This is part of a \$506,651 multi-institutional grant with UChicago, and co-PI Robert Rand; UMD is the lead.)

\item NSF CCF-1955610: Collaborative Research: SHF: Medium: Bringing Python Up to Speed, \$374,390, July 2020-2023. PI: Michael Hicks. (This is part of a \$1.2M multi-institution grant with UPenn and UMass Amherst; UMD is the lead institution.)

\item Advancing Checked C, Microsoft Research Award (gift), \$119,477. March 2019-December 2020. PI: Michael Hicks.

\item DARPA FA87501910008: Symbolic Input Unification and Minimization, \$34,802, Nov 2018 -- Dec 2018. PI: Michael Hicks.

\item NSF CNS-1801545 SaTC: CORE: Medium: Collaborative: Understanding Security in the Software Development Lifecycle: A Holistic, Mixed-Methods Approach, \$698,957, Sep 2018 -- 2022. Lead PI: Michelle Mazurek, co-PI: Michael Hicks. (This is part of a multi-institution grant with University of South Florida.)

\item DOE ASCR grant: Efficient and Reliable Mapping of Quantum Computations onto Realistic Architectures, \$4,467,131, Sep 2018 -- 2022. Lead PI: Andrew Childs. Co-PIs: Michael Hicks, Xiaodi Wu, Alexey Gorshkov.

\item Contesting Secure Development Effectiveness, Google Research Award (gift), \$55,256. March 2017--2018. PI: Michelle Mazurek, co-PI: Michael Hicks.

\item NSF CNS-1563722 TWC: Medium: Collaborative: New Protocols and Systems for RAM-Based Secure Computation,  \$464,196. May 2016--2019. Lead PI: Jonathan Katz, co-PI: Michael Hicks. (This is part of a multi-institution grant with George Mason University.)

\item DARPA FA87501520104: SOUCIS: Sound Over- \& Under-Approximations of Complexity \& Information Security, \$1,209,032. April 2015--2018. Lead PI: Michael Hicks, co-PI: David Van Horn and Jeffrey S. Foster. (This is part of a \$2.52M multi-institution grant with UC Berkeley and Yale, each with one additional co-PI; UMD is the lead institution.)

\item DARPA FA8750-16-C-0022: TAMBA: Testing and Modeling of Brandeis Artifacts, \$521,524. October 2015--October 2019. PI: Michael Hicks. (This is part of a multi-institution grant led by Galois, Inc.; Hicks is the only UMD PI.)

\item Connecting the Theory and Practice of Incremental Computation via Servo, gift from Mozilla Corporation, \$88,294. May 2014--2015. co-PIs: Michael Hicks and Jeff Foster.

\item Boombox: Dynamic Software Updates for Software Defined Networks, Google Research Award (gift), \$51,037. September 2014--2015. PI: Michael Hicks.

\item NSF CNS-1314857 TWC: Medium: Collaborative: DIORE: Digital Insertion and Observation Resistant Execution, \$799,499. August 2013-2018. PI: Elaine Shi, co-PIs: Michael Hicks and Bobby Bhattacharjee. (This is part of a \$1.2M multi-institution grant with one other PI at UT Austin. Bhattacharjee became UMD PI when Shi left UMD.)

\item NSF CNS-1111599 TC: Large: Collaborative Research: Practical Secure Two-Party Computation: Techniques, Tools, and Applications, \$1M. August 2011-2018. PI: Jonathan Katz, co-PI: Michael Hicks. (This is part of a multi-institution grant, with UVA as lead institution.)

\item Establishing a Science of Security Research Lablet at the University of Maryland, NSA, \$1,487,608. February 2014--February 2017. Lead PI: Jonathan Katz, co-PI: Michael Hicks (among 11 others).

\item Secure Information Flows in Hybrid Coalition Networks, US Army Research Lab and UK Ministry of Defence International Technology Alliance in Network and Information Science (USUKITA) program BPP15, \$244,233 PI: Michael Hicks, co-PI Jonathan Katz. May 2015--2016.

\item NSF EDU-1319147 EDU: Competing to Build Secure Systems, \$300,000. September 2013-2016. PI: Michael Hicks, co-PIs: Atif Memon, David M. Levin, and Jandelyn D. Plane.

\item Secure Information Flows in Hybrid Coalition Networks, US Army Research Lab and UK Ministry of Defence International Technology Alliance in Network and Information Science (USUKITA) program BPP13, \$356,615 PI: Michael Hicks, co-PI Jonathan Katz. May 2013--2015.

\item University of Maryland Partnership with the Laboratory of Telecommunications Sciences, Contract Number H9823013D00560002, Protecting against Malware on Android, \$371K. April 5, 2012--September 30, 2014. PI: Jeffrey S. Foster, co-PI: Michael Hicks.

\item NSF CCF-0910530 SHF: Large: Collaborative Research: Ever Ready: Perpetually Available Software Systems, \$642K.  August 2009--2014. PI: Michael Hicks. (This is part of a multi-institutional team grant with two other PIs for a total award of \$2.3M.)

\item NSF CCF-0915978 SHF: Small: User-Centered Software Analysis Tools, \$500K. September 2009-- 2013. PI: Jeffrey S. Foster, co-PI: Michael Hicks.

\item Securing Information Flows, US Army Research Lab and UK Ministry of Defence International Technology Alliance in Network and Information Science (USUKITA) program BPP11, \$333K. co-PIs: Michael Hicks and Jonathan Katz. May 2011--2013.

\item University of Maryland Partnership with the Laboratory for Telecommunications Sciences (LTS), Dynamic Software Update Automation for Servers and Event-Driven Programs, \$407K. January 2010-- 2013. PI: Michael Hicks, co-PI: Jeff Foster.

\item NSF CNS-0905419 RECOVERY: TC: Medium: Collaborative Research: Techniques to Retrofit Legacy Code with Security, \$300K. September 2009--2012. PI: Michael Hicks. (This is part of a multi-institutional team grant with three other PIs for a total award of \$1.2M.)

\item Improving the efficiency of static analysis of C source code using multiple tools, Cisco Systems Inc (Unrestricted gift), \$35,000. 2007. PI: Jeffrey S. Foster, co-PI: Michael Hicks.

\item NSF CCF-0541036 Scalable, Precise, and Effective Analyses for Detecting Race Conditions, \$360K. September 2006--2010. PI: Michael Hicks, co-PI: Jeffrey S. Foster.

\item US Army Research Lab and UK Ministry of Defence International Technology Alliance in Network and Information Science (USUKITA) program BPP09, \$197K. co-PIs: Michael Hicks and Jonathan Katz. May 2009--2011.

\item NSF CNS-0346989 CAREER: Programming Languages for Reliable and Secure Low-level Systems, \$550K. June 2004-2010. PI: Michael Hicks.

\item University of Maryland Partnership with the Laboratory for Telecommunications Sciences (LTS), Safe and Robust Dynamic Software Updating for Real-time Embedded Applications, \$327K, August 2007--January 2010. PI: Michael Hicks, co-PI: Jeff Foster.

\item Army Research Lab Communications and Networking Collaborative Technology Alliance (ARL C\&N CTA) program TIP-07-4, End-to-End Security of Cross-Domain Information Flows, \$245K (funds to Maryland). October 2006--October 2009. PI: Michael Hicks.

\item ARL C\&N CTA program (supported by DARPA DSO), Provably Enforcing Decentralized, Multi-level Security Policies in Distributed, Web-based Applications, \$83K. March 1, 2007--September 30, 2009. PI: Michael Hicks.

\item NSF CCF-0524036 Collaborative Research: CT-T: Flexible, Decentralized Information-flow Control for Dynamic Environments, \$280K. September 2005-2009. PI: Michael Hicks. (This is part of a multi-institutional team grant with three other PIs for a total award of \$1.1M.)

\item DARPA \#HR00110610019 Computer Science Study Group Panelist, \$77K. March 2006-2007. PI: Michael Hicks.

\item NSF IIS-0613601 SoD-HCER: Evaluation of Complex Designs--A Comparative Study, \$100K. July 2006--2008. PI: Michael Hicks (transferred from Vibha Sazawal in January 2007).

\end{description}


% \begin{description}
% \setlength{\labelsep}{0ex}
% \setlength{\itemsep}{-0.5ex}

% \item {\bf\textsf Collaborators} 

% \item { Thesis advisor}: Scott Nettles (Texas)

% \item { Postdoctoral advisor}: Greg Morrisett (Cornell)

% \item Former students and post-docs: 
%   Nikhil Swamy (Microsoft),
%   Iulian Neamtiu (UC Riverside),
%   Khoo Yit Phang (Mathworks),
%   Polyvios Pratikakis (FORTH, Greece),
%   Saurabh Srivastava (Berkeley),
%   Stephen Magill (Galois),
%   Edward Smith (UMass),
%   Suriya Subramanian (Microsoft),
%   Avik Chaudhuri (Facebook),
%   Nataliya Guts,
%   Matthew Hammer (CU Boulder),
%   Piotr Mardziel (Carnegie Mellon),
%   Karla Saur (Intel Labs),
%   Luis Pina (Imperial College, London)

% \item Other collaborators (2012--present): 
%   Yan Huang (Indiana),
%   Gavin Bierman (Oracle),
%   Daan Leijen (Microsoft),
%   Michael Clarkson (Cornell),
%   Mario Alvim,
%   Luis Veiga (IST, Portugal),
%   Emery Berger (UMass),
%   Gavin Bierman (Oracle Labs),
%   Nate Foster (Cornell),
%   Daan Leijen (Microsoft),
%   Kathryn McKinley (Microsoft/Texas),
%   Mudhakar Srivatsa (IBM),
%   Arjun Guha (UMass),
%   David Evans (UVa),
%   Laurent Vanbever (ETHZ),
%   Joshua Dunfield (UBC),
%   Elaine Shi (Cornell),
%   Mohit Tiwari (UT Austin)
%   Luis Veiga (IST Lisbon),
  
% \end{description}

\end{document}





